\section*{Resumen ejecutivo}
\addcontentsline{toc}{section}{Resumen ejecutivo}

{\sffamily\fontseries{b}\fontshape{n}\color{MASA-Blue} \textbf{Contexto y problema}}\\
El reconocimiento automático de matrículas vehiculares en México enfrenta tres barreras principales que han limitado su adopción. La primera es la ausencia de bases de datos públicas etiquetadas que respeten la diversidad normativa de los 32 estados, cuyos formatos incluyen guiones en posiciones variables y exclusiones de caracteres específicos como las letras \quotes{I}, \quotes{Ñ}, \quotes{O} y \quotes{Q}. La segunda se deriva de las restricciones éticas y legales asociadas a la captura y almacenamiento de imágenes reales de vehículos y sus ocupantes, lo que dificulta la recolección de datos de entrenamiento a gran escala. La tercera, y quizás la más relevante para el despliegue práctico, es que la mayoría de los sistemas de reconocimiento de matrículas existentes, tanto comerciales como de código abierto, han sido diseñados para ejecutarse en hardware de última generación (servidores con GPUs de alta capacidad o equipos de escritorio con procesadores potentes), lo que eleva significativamente los costos de infraestructura, mantenimiento y consumo energético, y los hace inviables para entornos con recursos limitados como municipios de bajos ingresos, estacionamientos residenciales, o aplicaciones embebidas distribuidas.

{\sffamily\fontseries{b}\fontshape{n}\color{MASA-Blue} \textbf{Solución propuesta}}\\
Nuestro sistema LPR-Sentinel, un sistema de reconocimiento de matrículas mexicanas, apuesta por operar exclusivamente con datos sintéticos generados conforme a la NOM-001-SCT-2-2016, eliminando así las complejidades éticas y legales del entrenamiento con imágenes reales. La arquitectura se compone de cinco módulos independientes: generador procedimental de matrículas por estado, aumentador de datos con 29 transformaciones, detector basado en YOLOv11s, reconocedor basado en FastPlateOCR, y un pipeline de integración con consultas a base de datos SQLite. Todos los modelos entrenados se exportan al formato de Intercambio de Redes Neuronales Abiertas (ONNX), permitiendo desacoplar el entrenamiento (realizado usando Tensorflow y PyTorch) de la inferencia (ejecutada sobre una Raspberry Pi 4B, un hardware embebido de bajo costo).

{\sffamily\fontseries{b}\fontshape{n}\color{MASA-Blue} \textbf{Resultados clave}}\\
La evaluación del sistema sobre 100 imágenes simuladas que incluyen condiciones extremas de perspectiva, iluminación y calidad arrojó una tasa de éxito extremo a extremo (detección correcta, reconocimiento exacto y consulta exitosa a la base de datos) del 56\% en escenas diurnas y del 36\% en escenas nocturnas, lo que representa una tasa global del 46\%. Estos valores, aunque deficientes en términos absolutos, deben interpretarse considerando dos factores relevantes. Primero, las imágenes de prueba incluyeron casos de muy alta dificultad como rotaciones superiores a 60 grados, perspectivas de captura extremas y oclusiones por sombras y luces parciales (la evaluación se realizó con toda intención de hacer fallar al sistema). Segundo, los modelos fueron entrenados con recursos computacionales limitados (solo 14 épocas para el detector y 115 épocas para el reconocedor), lo que sugiere que existe un amplio margen de mejora. De hecho, cuando se analiza el rendimiento de los módulos individuales en condiciones controladas, es decir, suprimiendo las imágenes donde el detector falla por completo por no detectar una matrícula, el detector alcanza una precisión excelente del 98.7\% y un mAP@0.5 de 0.992, mientras que el reconocedor logra una precisión por matrícula completa del 97.78\% y una precisión por carácter del 99.59\%, valores deseados dentro del estado del arte para tareas de reconocimiento óptico de caracteres en condiciones ideales.

En cuanto al rendimiento computacional, en un equipo de referencia (AMD Ryzen 7 5700G con 32 GB de RAM), la latencia media de 121 ms por imagen con un rendimiento de 8.26 imágenes por segundo son más que suficientes para aplicaciones de tiempo real, con un consumo moderado de 221\% de CPU (equivalente a 2.2 núcleos de los 8 disponibles para este modelo de procesador) y 276 MB de RAM por inferencia. En la plataforma objetivo, una Raspberry Pi 4B con 4 GB de RAM, el tiempo de procesamiento se incrementa a 1,428 ms por imagen, lo que representa un factor de ralentización de 11.8 veces y reduce el rendimiento a solo 0.70 imágenes por segundo (42 imágenes por minuto). Este rendimiento es insuficiente para aplicaciones que requieran respuesta instantánea, pero resulta aceptable para tareas de procesamiento asíncrono. Es importante destacar que la precisión del sistema se mantiene idéntica en ambas plataformas al utilizar los mismos modelos ONNX, lo que confirma que la degradación del rendimiento afecta exclusivamente a la velocidad, no a la calidad de las predicciones.

{\sffamily\fontseries{b}\fontshape{n}\color{MASA-Blue} \textbf{Innovación y cumplimiento ético}}\\
LPR-Sentinel constituye, hasta nuestro conocimiento, la primera implementación de código abierto que genera una base de datos sintética de matrículas mexicanas con estricto apego a la normativa nacional, entrenando modelos de detección y reconocimiento sin necesidad de capturar imágenes reales de ciudadanos. Esta propuesta ha demostrado ser funcional, alcanzando una precisión del 98.7\% en detección y del 97.78\% en reconocimiento por matrícula completa cuando se evalúan en condiciones controladas. Estos resultados respaldan nuestra hipótesis en la que buscamos demostrar que los datos sintéticos, cuando son generados con estricto apego a una normativa y aumentados adecuadamente, pueden ser una alternativa viable y éticamente superior a los datos reales para el entrenamiento de sistemas de reconocimiento de matrículas, incluso con la heterogeneidad presente en México.

{\sffamily\fontseries{b}\fontshape{n}\color{MASA-Blue} \textbf{Viabilidad de despliegue}}\\
El sistema ha sido validado en una Raspberry Pi 4B con 4 GB de RAM y almacenamiento SSD por USB 3.0, operando con un consumo energético inferior a 5 watts. Para contextualizar el ahorro económico que esto representa, un servidor tradicional equipado con una GPU NVIDIA T4 (configuración típica en el ecosistema de Google Cloud) tiene un costo mensual de aproximadamente 150 dólares y consume entre 150 y 300 watts en operación continua. En cambio, nuestro hardware objetivo (Raspberry Pi 4B, unidad de almacenamiento y sistema de instalación) tiene un costo inferior a 250 dólares (único pago) y consume en promedio 5 watts, lo que representa un único costo de adquisición en lugar de la renta de hardware y se reduce hasta en un 97\% en consumo energético.

Aunque la latencia medida en la Raspberry Pi 4B (1,428 ms por imagen, equivalente a 42 imágenes por minuto) es insuficiente para aplicaciones que requieren respuesta instantánea, sigue resultando aceptable para múltiples casos de uso del mundo real que no demandan latencia crítica. Entre ellos se incluyen el control de acceso a estacionamientos residenciales de bajo flujo vehicular, la revisión de entrada y salida de vehículos en estaciones de registro o en sistemas de seguridad para comunidades rurales o de bajos recursos donde la inversión en hardware especializado no es viable. Adicionalmente, el sistema una vez adquirido permite futuras optimizaciones que podrían reducir la latencia sin cambiar la plataforma hardware. Como alternativa de mediano plazo, el sistema puede migrar sin modificación alguna de código a plataformas más capaces dentro del mismo rango de costo, como la Raspberry Pi 5 (estimado de 400-500 ms por imagen), la Orange Pi 5 (estimado de 80-120 ms por imagen), o la NVIDIA Jetson Nano (estimado de 50-100 ms por imagen), todas ellas compatibles con los mismos modelos ONNX generados en nuestra propuesta.

{\sffamily\fontseries{b}\fontshape{n}\color{MASA-Blue} \textbf{Próximos pasos}}\\
El análisis de los resultados obtenidos ha permitido identificar las principales deficiencias del sistema actual. El módulo de detección basado en YOLOv11s, si bien alcanza una precisión excelente del 98.7\% en condiciones controladas, es responsable de la mayor parte del consumo de recursos durante la inferencia: entre 800 y 1,200 milisegundos por imagen en la Raspberry Pi 4B y más de 100 megabytes de memoria RAM por ejecución. Este desbalance es la causa principal de que la latencia total del sistema alcance los 1,428 milisegundos, limitando su aplicabilidad en tiempo real. Por lo que, como siguiente paso inmediato, se planea reemplazar YOLOv11s por el modelo Vision Transformer RF-DETR, una arquitectura más moderna que ha sido reportada en la literatura como significativamente más eficiente que las variantes YOLO en tareas de detección de objetos, especialmente en dispositivos con recursos computacionales limitados. Se espera que este reemplazo, sin necesidad de modificar el resto del sistema, reduzca la latencia del detector a un rango de 200 a 400 milisegundos en la Raspberry Pi 4B, lo que acercaría el tiempo total de procesamiento a valores compatibles con aplicaciones de tiempo real poco exigentes (3 a 5 imágenes por segundo).

Paralelamente, se llevará a cabo un reentrenamiento completo de ambos modelos (detector y reconocedor) incorporando las mejoras necesarias para abordar las deficiencias identificadas durante la primera evaluación. Estas incluyen la inclusión de transformaciones de aumentación más agresivas para simular rotaciones extremas (hasta 90 grados), desenfoques de movimiento más intensos y condiciones de iluminación más intensas, así como un aumento en el número de épocas de entrenamiento (de 14 a 100 o más para el detector, y de 115 a 300 para el reconocedor), aprovechando recursos computacionales más amplios que los disponibles durante la fase inicial de desarrollo. Se espera que estas mejoras eleven la tasa de éxito en condiciones nocturnas desde el 36\% actual hasta al menos el 55-60\%, acercándose al rendimiento diurno y reduciendo la brecha entre ambos escenarios.

A mediano plazo, se planea validar el sistema con imágenes reales capturadas en colaboración con autoridades de tránsito y cuerpos de seguridad (al menos con el estado de Guanajuato, estado actual en el que se radica), bajo los protocolos éticos correspondientes que garanticen el anonimato y la protección de datos. Esta validación permitirá comparar cuantitativamente el rendimiento de los modelos entrenados exclusivamente con datos sintéticos frente a su desempeño en el mundo real. Adicionalmente, se incorporará un módulo de seguimiento temporal que permita asociar detecciones de la misma placa en secuencias de video, eliminando la actual restricción de operar con imágenes estáticas ocasionada por la actual configuración de optimización del dispositivo empleado. Finalmente, se extenderá el generador sintético para soportar otros tipos de vehículos actualmente circulantes en México, incluyendo servicio público (taxis, transporte de personal), vehículos oficiales de gobierno, unidades para personas con discapacidad, y transporte de carga pesada, cada uno con sus propias combinaciones de colores, formatos y rangos alfanuméricos, sin modificar la arquitectura general del sistema.

\vspace{1cm}
{\sffamily\fontseries{b}\fontshape{n}\color{MASA-Blue} \textbf{Contacto del desarrollador:}} \\
Gustavo Adolfo Murillo Gutierrez\\
ga.murillogutierrez@outlook.com\\
(+52) 462 400 2500 \\
Irapuato, Guanajuato, México. \\
GitHub personal: \url{https://github.com/MurilloLog/} \\
Video de demostración: \url{[enlace]}